
\begin{thebibliography}{99}

\bibitem{WHO2021} 世界卫生组织. 全球卫生统计报告2021[R]. 日内瓦：WHO出版社, 2021.

\bibitem{zheng2023} 郑月华, 黄新华, 李建军. 慢性病防控信息化建设的进展与展望[J]. 中国医学文摘·内科学, 2023, 45(3): 234-241.

\bibitem{liu2022} 刘晓芳, 王林, 张静. 基于机器学习的慢性病风险预测模型研究进展[J]. 医学信息学杂志, 2022, 43(5): 12-19.

\bibitem{chen2023} 陈卫红, 李强, 王美琳, 等. 患者健康数据标准化处理与质量管理体系[J]. 中华流行病学杂志, 2023, 44(2): 187-195.

\bibitem{patel2021} Patel M S, Volpp K G, Asch D A. Wearable devices as facilitators, not drivers, of health behavior change[J]. JAMA, 2021, 313(5): 459-460.

\bibitem{Reges2015} Reges O, Enav T, Sherf M, et al. Association between net user rates and the incidence of metabolic syndrome: A nationwide study[J]. JAMA Network Open, 2015, 318(18): 1857-1858.

\bibitem{rasheed2021} Rasheed K, Qayyum A, Qadir J, et al. Machine learning for Internet of Things: A survey[J]. Internet of Things, 2021, 3-4: 100245.

\bibitem{wang2023} 王志红, 吴雪, 王琳. 个性化健康管理方案的设计与实施[J]. 中华健康管理学杂志, 2023, 17(4): 298-305.

\bibitem{Li2022} Li J, Wang H, He Y, et al. Data visualization for health information systems: A systematic review[J]. Healthcare Informatics Research, 2022, 28(1): 3-15.

\bibitem{guo2023} 郭春燕, 张美芳, 李红梅. 基层医疗机构慢性病管理现状与改进对策[J]. 中国全科医学, 2023, 26(8): 945-952.

\bibitem{friedman2000} Friedman J H. Greedy function approximation: A gradient boosting machine[J]. Annals of Statistics, 2000, 29(5): 1189-1232.

\bibitem{breiman2001} Breiman L. Random forests[J]. Machine Learning, 2001, 45(1): 5-32.

\bibitem{GDPR2018} 欧盟. 一般数据保护条例(GDPR)[S]. 欧盟官方公报, 2018.

\bibitem{HIPAA1996} 美国. 健康保险可携性和问责法(HIPAA)[S]. 美国联邦法规, 1996.

\bibitem{zhu2022} 朱长久, 李文娟, 陈建华. Python在医疗大数据分析中的应用[J]. 中国医学装备, 2022, 19(5): 45-51.

\end{thebibliography}
